Dieses Experiment untersucht die Winkelverteilung und die mittlere Lebensdauer von Myonen in der kosmischen Strahlung. Zusätzlich wird das Amplitudenspektrum der detektierten Myonenpulse aufgezeichnet, um deren Energieverteilung und die Eigenschaften des Detektionssystems zu analysieren.

Ein Schlüsselelement des Experiments ist die präzise Abstimmung der Diskriminatoren, einschließlich der Bestimmung der optimalen Schwellenspannungen zur Unterdrückung von Rauschspektren und Störereignissen. Des Weiteren werden zufällige Koinzidenzen quantifiziert, um die tatsächliche Myonengeschwindigkeit präzise von zufälligen zeitlichen Überlappungen zu unterscheiden und die Zuverlässigkeit der Messdaten zu bewerten.